% Programming with InitechBase, Chapter 16: Working with Dates (AN-16).
% Panels replayed (examples/p16-*.dot); sources core/eval.c (date - date,
% date + n), core/fn_builtins.c (CTOD, DOW, CDOW, DAY), the program
% SERVICE.PRG.
\chapter{Working with Dates}
\Gates{test-pubs-examples, test-xbase-fn-b, test-xbase-eval}
Personnel records are full of dates: of hire, of review, of birth.  This
chapter computes the years of service of every employee, as of a date the
program is given.

\begin{ChapterCovers}
\item Date arithmetic
\item The day of the week
\item Years of service
\end{ChapterCovers}

\ATEHead{Preparing for This Chapter}
Start InitechBase with a fresh copy of the sample diskette, and type
\type{USE STAFF}.

\ATEHead{Date Arithmetic}
One date minus another is the number of days between them; a date plus or
minus a number of days is a date.  The calendar is followed exactly, leap years
included:
\DotExample{examples/p16-dates}
\fn{DOW()} numbers the days of the week from 1, Sunday, to 7, Saturday;
\fn{CDOW()} names them.

\ATEHead{Years of Service}
\file{SERVICE.PRG} computes, for each employee, the complete years of service
as of the date in \texttt{MASOF}, and counts the staff in three classes:
\ProgramListing{SERVICE.PRG}
\begin{Steps}
\item Type \type{MASOF = CTOD('03/01/92')}.
\item Type \type{DO SERVICE}.
\end{Steps}
\DotExample{examples/p16-service}
The days of service divided by 365, with the fraction dropped by \fn{INT()},
give the years.  Like the aging report of Chapter~9, the program takes its date
from a memory variable, so it can be run as of any date.

\ATEHead{Where Do You Go From Here?}
Chapter~17 works with the characters of names.

\ATEHead{Summary}
\begin{enumerate}
\item Date minus date gives days; date plus days gives a date.
\item \fn{DOW()} and \fn{CDOW()} give the day of the week.
\item \fn{INT()} of the days divided by 365 gives complete years.
\end{enumerate}
